\chapter{Kitab Thaharah}

\section{Definisi Thaharah}

\textbf{Makna kitab thaharah:} Kumpulan hukum-hukum thaharah.

\textbf{Thaharah secara syar'i:} Mengangkat hadats atau menghilangkan najis, dan apa-apa yang semakna dengan keduanya atau yang serupa dengan keduanya.

\textbf{Najis:} Suatu kotoran yang menjijikkan menghalangi sahnya shalat, selama tidak ada rukhshah; yaitu keringanan yang membolehkannya, seperti kotoran tinja manusia itu sendiri.

\textbf{Mutananjis:} Sesuatu yang suci yang terkena najis seperti pakaian dan karpet.

\textbf{Mengangkat hadats:} Menghilangkan sifat syar'i yang melekat pada badan yang menghalangi sahnya shalat atau semisalnya. Contohnya: Wudhu untuk mengangkat hadats kecil dan mandi untuk mengangkat hadats besar. 

Hadats dibagi menjadi 2 jenis:
\begin{enumerate} 
    \item \textit{Hadats asghar} (hadats kecil): yaitu hadats yang diangkat wudhu.
    \item \textit{Hadats akbar} (hadats besar): yaitu hadats yang diangkat dengan mandi.
\end{enumerate}

Maksud dari \textit{"dan apa-apa yang semakna dengan keduanya"}; yaitu apa-apa yang semakna dengan mengangkat hadats dan yang semakna dengan menghilangkan najis.

\begin{enumerate}
    \item Apa-apa yang semakna dengan mengangkat hadats:
    \begin{itemize}
        \item Memandikan jenazah. Karena memandikannya tidak disebabkan oleh hadats, tetapi sebagai perkara \textit{taabud} (sebagai bentuk ibadah) dan perbuatan itu tidak mengangkat hadats.
        \item Tayamum. Dalam madzhab, tayamum tidak mengangkat hadats, melainkan hanya membolehkan shalat dan yang semisalnya yang diharamkan sebab hadats. Oleh karena itu mereka mengatakan: Tayamum bersifat \textit{mubih} (membolehkan) dan bukan mengangkat hadats. Sedangkan dalam Mukhtasar Abi Qatadah Fi Fiqhi Asy-Syafi’i As-Sadah: Tayamum seperti wudhu yang mengangkat hadats.
    \end{itemize} 
    \item Apa-apa yang semakna dengan menghilangkan najis: Seperti mengusap, menggosok dan dengan sinar matahari.
\end{enumerate}
	
Maksud dari \textit{yang serupa dengan keduanya}; yaitu serupa dalam mengangkat hadats dan serupa dalam menghilangkan najis:
\begin{enumerate}
    \item Yang serupa dalam mengangkat hadats:
    \begin{itemize}
        \item Basuhan kedua, dan ketiga dalam berwudhu. Hadats sebenarnya telah terangkat pada basuhan pertama. Adapun basuhan kedua dan ketiga adalah sunah inilah yang disebut \textit{yang serupa dengan mengangkat hadats} yang pada hakikatnya tidak mengangkat hadats.
        \item Memperbarui wudhu.
        \item Mandi yang disunahkan seperti mandi jum’at.
    \end{itemize}
    \item Yang serupa dalam menghilangkan najis: Seperti basuhan kedua dan ketiga pada tempat yang mengeluarkan najis (dubur dan qubul) setelah kita menghilangkan zat najisnya.
\end{enumerate}

\textbf{Hukum Thaharah:} Bersuci dari najis dan menghilangkannya hukumnya wajib bagi yang ingat dan mampu. Adapun bersuci dari hadats hukumnya wajib untuk membolehkan shalat.

\section{Tujuan, Sarana, dan Perantara Sarana Thaharah}

Tujuan thaharah ada empat: 
\begin{enumerate}
    \item Wudhu.
    \item Mandi. 
    \item Tayamum. 
    \item Menghilangkan najis.
\end{enumerate}

Sarana thaharah ada empat: 
\begin{enumerate}
    \item Air.
    \item Debu.
    \item Batu. 
    \item Menyamak kulit.
\end{enumerate}

Perantara sarana thaharah ada dua: 
\begin{enumerate}
    \item Bejana.
    \item Ijitihad.
\end{enumerate}
